\section{Retained barriers and impossibility results}
\label{sec:barriers}

A general framework is credible only if it states what encoding cannot repair.
The following barriers explain why the positive interfaces in earlier sections
are deliberately guarded, relational, or certificate-relative.

\subsection{No total universal primitive-recursive evaluator}

\begin{theorem}[Diagonal barrier]
\label{thm:no-universal-pr}
There is no binary primitive-recursive function \(U(e,x)\) such that every
unary primitive-recursive function \(f\) has an index \(e\) with
\(f(x)=U(e,x)\) for all \(x\).
\end{theorem}

\begin{proof}
If such \(U\) existed, then
\[
 g(x)=U(x,x)+1
\]
would be primitive recursive.  By universality, \(g(x)=U(e,x)\) for some
\(e\).  Substitution of \(x=e\) gives
\[
 U(e,e)+1=g(e)=U(e,e),
\]
a contradiction.
\end{proof}

This does not prevent a syntax of primitive-recursive programs with a
partial, clocked, or complexity-bounded interpreter.  It prevents one fixed
total primitive-recursive evaluator from uniformly evaluating every such
program.  Therefore the framework uses fixed evaluator schemas or finite
traces whose individual steps are checked.

\subsection{A cardinality bound is not a search cutoff}

Let \(T(e,t)\) be the primitive-recursive predicate saying that machine \(e\)
halts for the first time at exactly stage \(t\).  For each \(e\), at most one
\(t\) satisfies \(T(e,t)\).  Thus the solution set has the uniform cardinality
bound one.

\begin{theorem}[No uniform cutoff from cardinality]
\label{thm:cardinality-no-cutoff}
There is no total computable function \(H(e)\) such that
\[
 T(e,t)\to t\le H(e)
\]
for every \(e,t\).
\end{theorem}

\begin{proof}
Given such \(H\), simulate machine \(e\) for \(H(e)\) stages.  If it halts,
answer yes.  If it has not halted, it can never halt later, because its first
halting stage would exceed the promised bound.  This decides the halting
problem.
\end{proof}

Consequently, a theorem proving that a curve has at most \(N\) rational points
does not give a height bound, an enumeration stopping rule, or the points
themselves.  Additional arithmetic structure may produce a height bound, but
cardinality alone cannot.

\subsection{Exact sign and LLPO}

Suppose a total operation on arbitrary fast Cauchy names returned one of the
two valid conclusions
\[
 x\le0\qquad\text{or}\qquad x\ge0.
\]
Given a binary sequence \(\alpha\) with at most one nonzero term, form the fast
Cauchy real
\[
 x_\alpha=\sum_{n\ge0}(-1)^n\alpha_n2^{-n-1}.
\]
If the operation returns \(x_\alpha\ge0\), no odd-indexed term can be one; if
it returns \(x_\alpha\le0\), no even-indexed term can be one.  At the zero
sequence both conclusions are valid, so either answer is acceptable.  This is
exactly the lesser limited principle of omniscience (LLPO).

Thus an exact side, sign, orientation, or branch selector at a zero boundary
is not obtained merely by taking more approximations.  A positive gap makes
the finite comparison valid; an overlap makes either branch sound; otherwise
the output should remain relational.

The rotation-axis obstruction is an instance.  Near the identity, the
oriented vector part of a quaternion can approach zero from incompatible
directions.  Exact unit-axis selection would decide a discontinuous side
problem.  Quaternion data with a possibly zero vector part, or approximate
axes, remain stable.

\subsection{Arbitrary Cauchy coefficients do not yield exact Puiseux data}

Consider, for a nonnegative represented real \(a\),
\[
 P_a(X,Y)=Y^2-X^2-a.
\]
Suppose a uniform algorithm returned a ramification exponent \(e>0\), a
finite Laurent lower bound, and exact computable names of coefficients of a
formal root
\[
 y(t)=\sum_j y_jt^j,
 \qquad y(t)^2=t^{2e}+a.
\]
Finite coefficient comparison gives the following dichotomy.

If \(a=0\), all coefficients below degree \(e\) vanish and
\(y_e^2=1\), so \(|y_e|=1\).  If \(a>0\), then \(y_0^2=a\) is apart from zero;
induction through degrees below \(2e\) forces \(y_e=0\).  An approximation to
the algorithm's own coefficient \(y_e\) therefore decides whether \(a=0\).

Now let \(a_m\) have the fixed primitive-recursive fast Cauchy approximation
\[
 A(m,s)=\sum_{k<s}2^{-k-1}H(m,k),
\]
where \(H(m,k)\) says that machine \(m\) has halted by stage \(k\).  If the
machine never halts, \(a_m=0\); if its first halting stage is \(h\), then
\(a_m=2^{-h}>0\).  The hypothetical Puiseux producer would decide halting.

The obstruction concerns exact individually accessible branch coefficients
from arbitrary Cauchy input.  It does not rule out approximate root clusters,
conditional simple branches, or exact Puiseux computation from finite
algebraic presentations supplying equality and splitting data.

\subsection{Metastability does not supply a Cauchy modulus}

There are primitive-recursive increasing bounded rational sequences with no
computable Cauchy modulus.  One standard construction uses stage-bounded
halting information:
\[
 s_n=\sum_{e<n\,:\,e\text{ halts within }n\text{ stages}}2^{-e-1}.
\]
The map \(n\mapsto s_n\) is primitive recursive, rational, increasing, and
bounded by one.  If a computable modulus \(M(r)\) made every later variation
smaller than \(2^{-r}\), then to decide whether machine \(e\) ever halts one
would inspect stage \(M(e+2)\).  A later first contribution of weight
\(2^{-e-1}\) would exceed the allowed future variation; hence absence by that
stage would imply nonhalting.  This is impossible.

Finite stable-window or metastability theorems can still be primitive
recursive and useful.  What fails is the uniform conversion of such windows
into one global convergence rate.  The shrinking-viability theorem therefore
requires a tail modulus as an explicit input.

\subsection{Qualitative uniqueness is not a quantitative selector}

A statement that exactly one mathematical object satisfies a predicate does
not automatically give a primitive-recursive rate at which finite candidate
clouds shrink.  The unique object may be encoded by a computation whose
convergence has no computable modulus.  To construct a name, the framework
requires an isolation radius, strong-convexity constant, residual-to-distance
bound, contraction factor, or equivalent quantitative determinacy data.

This is why the report distinguishes:

\begin{enumerate}
\item uniqueness as a logical proposition;
\item finite near-uniqueness at a requested scale; and
\item a primitive-recursive modulus witnessing how the scale improves.
\end{enumerate}

Only the third supports a uniform coherent-name constructor.

\subsection{A semantic Lean proof does not automatically reflect into iPRA}

Lean can prove an ambient theorem about a recursive function using induction,
quotients, classical choice, algebraic hierarchies, or library facts.  There is
no general reflection principle turning an arbitrary such theorem into a term
of \(\Prf\,\PRAAx\).  A soundness theorem goes in the other direction: an
object-language proof has a semantic interpretation.

The positive solution is deliberately smaller.  For the restricted DSL of
Boolean folds, bounded traces, and certified arithmetic transducers, define a
syntax-directed compiler that constructs the object-language proof from local
object-language obligations.  This is feasible because every proof step and
induction invariant is exposed.  It is not a universal reflection theorem for
Lean.

\subsection{Boundary summary}

The framework does not claim:

\begin{itemize}
\item a universal total primitive-recursive evaluator;
\item exact discontinuous selectors from arbitrary Cauchy names;
\item height cutoffs from cardinality bounds;
\item a Cauchy rate from metastability alone;
\item quantitative isolation from qualitative uniqueness alone; or
\item automatic reflection of arbitrary semantic Lean proofs.
\end{itemize}

Each refused claim has a positive replacement: fixed evaluators or trace
certificates, relational output or positive gaps, explicit height producers,
finite-window theorems or tail moduli, shrinking certificates, and a restricted
proof-producing compiler.
